\documentclass[10pt,a4paper]{amsart}
\usepackage[margin=19mm]{geometry}
\usepackage{amsmath,amssymb,mathtools}
\usepackage[scr=rsfs,cal=euler]{mathalfa}
\usepackage{xfrac}
\usepackage[unicode,hidelinks,bookmarks=false]{hyperref}
\newcommand{\ie}{i.e.\ }
\newcommand{\cf}{cf.\ }
\newcommand{\resp}{respectively\ }
\newcommand{\Subsection}{\subsection}
\providecommand{\nobreakdash}{\nobreak\hbox{-}}
\setlength{\emergencystretch}{2em}
\multlinegap0pt
\usepackage{bm}
\usepackage{bbm}
\usepackage{dsfont}
\usepackage{xspace}
\usepackage{cancel}
\usepackage{mathtools}
\usepackage{cleveref}
\usepackage{amsthm}
\makeatletter
\@ifundefined{theorem}{\newtheorem{theorem}{Theorem}}{}
\@ifundefined{assumption}{\newtheorem{assumption}[theorem]{Assumption}}{}
\makeatother
\makeatletter
\newcommand{\basemeasure}{\pi_{\text{knw}}}
\newcommand{\dataerror}{\mathsf{Red}_D}
\expandafter\ifx\csname ssubproof\endcsname\relax
\newenvironment{ssubproof}[1][\proofname]{%
  \renewcommand{\qedsymbol}{$\blacksquare$}%
  \begin{proof}[#1]%
}{%
  \end{proof}%
}
\fi
\expandafter\ifx\csname subproof\endcsname\relax
\newenvironment{subproof}[1][\proofname]{%
  \renewcommand{\qedsymbol}{$\square$}%
  \begin{proof}[#1]%
}{%
  \end{proof}%
}
\fi
\newcommand{\testdataerror}{\mathsf{Red}_D}
\makeatother
\title[Understanding LLM Behaviors via Compression: Data Generation]{Understanding LLM Behaviors via Compression: Data Generation, Knowledge Acquisition and Scaling Laws}
\author{Zhixuan Pan, Shaowen Wang, Jian Li}
\date{Source: arXiv:2504.09597v6 (CC BY 4.0)}
\begin{document}
\maketitle

\noindent\emph{Source and licence.} Understanding LLM Behaviors via Compression: Data Generation, Knowledge Acquisition and Scaling Laws. Version arXiv:2504.09597v6 (CC BY 4.0), distributed under the Creative Commons Attribution 4.0 International licence, as recorded in the arXiv licence metadata for this exact version and shown on the abstract page https://arxiv.org/abs/2504.09597v6. Full rights notice: licenses/arxiv-2504.09597-CC-BY-4.0.txt.\par\smallskip

\noindent\textit{Editorial scope.} Counted unit: the Theorem whose source label is \texttt{thm: data scaling lower bound} in the article (source0.tex lines 2605--2615, source label \texttt{thm: data scaling lower bound}), delivered together with the article's own complete original proof of it (source0.tex lines 2616--2693, one \texttt{proof} environment of 78 lines). No mathematics is written, repaired or re-derived: statement and proof are the article's own text line for line. Required context retained uncounted: ass:data law (lines 2572--2580). Every definition, display and label of those lines is retained unchanged. Omitted from this package and not redistributed: 1--2571, 2581--2604, 2694--4206. Nothing inside a retained excerpt is omitted or altered: every retained body is delivered exactly as the article's lines are written, so no omission receipt of any kind accompanies this package. Citations: the retained excerpts carry no citation command at all, so no bibliography entry of the article is transcribed and no reference is redistributed. Reference adaptation: none was needed and none was made; every reference inside the delivered unit resolves inside it, so no unresolved reference or citation is produced. Printed slips kept literal: no printed word, symbol, hypothesis or number of the counted statement or of its proof was altered. Layout and class: the article's own class is \texttt{article} with packages including setspace, url, booktabs, amsfonts, nicefrac, microtype, xcolor, amsmath, amsthm, graphicx, subcaption, bbm, natbib, hyperref; the wrapper is the lane's pinned \texttt{amsart} editorial wrapper at 10pt on A4, which restates the theorem-like environments the retained text needs and adds the article's own macro definitions that the retained text needs (\texttt{\textbackslash basemeasure}, \texttt{\textbackslash dataerror}, \texttt{\textbackslash testdataerror}), with extra wrapper packages (bm, bbm, dsfont, xspace, cancel, mathtools, cleveref). The delivered numbering is the unit's own. Assets: no retained span contains a figure environment, a graphics-inclusion command, a PDF-page inclusion, a table or a TikZ picture; no image file is copied into this package, the package declares no asset dependency and no image file is redistributed with it. Licence: version arXiv:2504.09597v6 of this article is distributed under the Creative Commons Attribution 4.0 International licence, as recorded in the arXiv licence metadata for this exact version and shown on the abstract page https://arxiv.org/abs/2504.09597v6. A full rights notice is delivered at licenses/arxiv-2504.09597-CC-BY-4.0.txt. Characters: every retained excerpt is pure ASCII, so the delivered engine, XeLaTeX with its default Latin Modern text font, reports no missing character.
\begin{assumption}
\label{ass:data law}
(\uppercase\expandafter{\romannumeral 1})
The expected redundancy of unseen cluster of knowledge (i.e., $t_i=0$) exceeds a certain constant $c_1$, i.e., $\mathbb{E}_{\phi_i\sim \basemeasure}\left[\dataerror^{(i)}(0)\right]>c_1$ for all $k\in \mathcal{N}^+$.

(\uppercase\expandafter{\romannumeral 2})
For all cluster of knowledge, there exists a constant $c_2$ such that 
 $\mathbb{E}_{\phi_i\sim \basemeasure}\left[\dataerror^{(i)}(t)\right]\le \dfrac{c_2}{t+1}$ for all $i,t\in \mathcal{N}^+$.
\end{assumption}

\begin{theorem}
\label{thm: data scaling lower bound}
Under (\uppercase\expandafter{\romannumeral 1}) in \Cref{ass:data law}, the total expected redundancy $\testdataerror$ satisfies
    \begin{align*}
    \testdataerror(N) = \Omega(N^{-1+\alpha}).
\end{align*}
Under (\uppercase\expandafter{\romannumeral 2}) in \Cref{ass:data law}, the total expected redundancy $\testdataerror$ satisfies
    \begin{align*}
    \testdataerror(N) = O(N^{-1+\alpha}).
\end{align*}
\end{theorem}

\begin{proof}[Proof of \Cref{thm: data scaling lower bound}]
% By \Cref{thm:PYP rank fre}, we know that there exists $N_0$ such that for all $n\ge N_0$, 
% $\dfrac{S_{\alpha,\beta}}{2n^{1/\alpha}}\le p_n\le \dfrac{2S_{\alpha,\beta}}{n^{1/\alpha}}$ almost surely. Then we have
Denote $f_i(t_i)=\mathbb{E}_{\phi_i\sim \basemeasure}[\dataerror^{(i)}(t_i)]$.

We begin by proving the first part of the theorem.
    \begin{align*}
        \testdataerror(N)&=\sum\limits_{k=1}^{\infty} p_k \mathbb{E}_{t_k}[f_k(t_k)]
        \\&\ge \sum\limits_{k> N^{\alpha}} p_k\mathbb{E}_{t_k}[f_k(t_k)] 
        \\& \overset{(1)}{=} \sum\limits_{k> N^{\alpha}}\sum\limits_{l=0}^N p_k^{l+1} (1-p_k)^{N-l}\binom{N}{l}f_k(l) \\
        &\overset{(2)}{\ge} \sum\limits_{k> N^{\alpha}} p_k (1-p_k)^{N} f_k(0)\\
        &\overset{(3)}{\ge} \dfrac{c_1}{\zeta(1/\alpha)}\sum\limits_{k> N^{\alpha}} \dfrac{1}{k^{1/\alpha}} \left(1-\dfrac{1}{\zeta(1/\alpha)N}\right)^{N} \\
        &\ge \dfrac{c_1}{2e\zeta(1/\alpha)}\sum\limits_{k> N^{\alpha}} \dfrac{1}{k^{1/\alpha}} \\
        &=  \Omega(N^{-1+\alpha}).
    \end{align*}
(1) is obtained by expanding the expectation; (2) corresponds to taking the term with $l = 0$; and (3) is derived by substituting the value of $p_k=\dfrac{1}{k^{1/\alpha}\zeta(1/\alpha)}$.

% \end{proof}


% \begin{theorem}[Restatement of \Cref{thm: data scaling upper bound}]
% Assume there exists a constant $C$ such that 
%  $(t+1)\dataerror(k,t)\le C$ for all $k,t\in \mathcal{N}^+$. Then the test error $\testdataerror$ satisfies
%     \begin{align*}
%     \testdataerror(N)\le O(N^{-1+\alpha}).
% \end{align*}
% \end{theorem}
% \begin{proof}[Proof of \Cref{thm: data scaling upper bound}]
% By \Cref{thm:PYP rank fre}, we know that there exists $N_0$ such that for all $n\ge N_0$, 
    
% $\dfrac{S_{\alpha,\beta}}{2n^{1/\alpha}}\le p_n\le \dfrac{2S_{\alpha,\beta}}{n^{1/\alpha}}$ almost surely.

    Next, we prove the second part. For all $k\in \mathcal{N}^+$, we know that 
    % \begin{align*}
    %     \sum\limits_{k\le (2S_{\alpha,\beta}N)^{\alpha}}\mathbb{E}_{\knwModel}\left[p_kE_{t_k}[\dataerror(k,t_k)]\right]  = O(N^{-1+\alpha}).
    % \end{align*}

    
    % Let the training data be $X$, and the model trained on $X$ has a test loss for knowledge $Y$ denoted by $g_X(Y)$.
    % \begin{align*}
    %     E_{X,Y}[g_X(Y)] &= E_{X}\left[\sum\limits_{k=1}^{\infty}p_kf(t_k)\right] \\
    %     &= \sum\limits_{k=1}^{\infty}p_kE_{X}[f(t_k)].
    % \end{align*}
    % Here, $t_k$ is the number of times the $k$-th knowledge appears in the training data.
    \begin{align*}
        p_k\mathbb{E}_{t_k}[f_k(t_k)] &\overset{(1)}{=} p_k\sum\limits_{l=0}^N p_k^l (1-p_k)^{N-l}\binom{N}{l}f_k(l) \\
        &= \sum\limits_{l=0}^N p_k^{l+1} (1-p_k)^{N-l}\dfrac{l+1}{N+1}\binom{N+1}{l+1}f_k(l) \\
        &\le \dfrac{\max\limits_{l} (l+1)f_k(l)}{N+1}\sum\limits_{l=0}^N p_k^{l+1} (1-p_k)^{N-l}\binom{N+1}{l+1} \\
        &\overset{(2)}{\le} \dfrac{\max\limits_{l} (l+1)f_k(l)}{N+1}\\ &\le \dfrac{c_2}{N+1}.
    \end{align*}
    (1) is obtained by expanding the expectation, and (2) is derived by utilizing 
    \\
    $\sum\limits_{l=0}^N p_k^{l+1} (1-p_k)^{N-l}\binom{N+1}{l+1}\le \sum\limits_{l=-1}^N p_k^{l+1} (1-p_k)^{N-l}\binom{N+1}{l+1}=1.$

    Therefore,
    \begin{align*}
        \sum\limits_{k\le N^{\alpha}}p_k\mathbb{E}_{t_k}[f_k(t_k)] \le  \dfrac{c_2}{N+1} N^{\alpha} =O(N^{-1+\alpha}).
    \end{align*}
    For $k>N^{\alpha}$, we have 
    \begin{align*}
        \sum\limits_{k> N^{\alpha}}p_k\mathbb{E}_{t_k}[f_k(t_k)] &= \sum\limits_{k> N^{\alpha}}\sum\limits_{l=0}^N p_k^{l+1} (1-p_k)^{N-l}\binom{N}{l}f_k(l) \\
        &\overset{(1)}{\le}\ \sum\limits_{k> N^{\alpha}}\sum\limits_{l=0}^N \left(\dfrac{1}{\zeta(1/\alpha)k^{1/\alpha}}\right)^{l+1} \binom{N}{l}f_k(l) \\
        &\overset{(2)}{\le} \sum\limits_{k> N^{\alpha}}\sum\limits_{l=0}^N \left(\dfrac{1}{k^{1/\alpha}}\right)\left(\dfrac{1}{N}\right)^l \binom{N}{l}f_k(l) \\
        &\overset{(3)}{\le}\ \sum\limits_{k> N^{\alpha}}\sum\limits_{l=0}^N \left(\dfrac{1}{k^{1/\alpha}}\right)\dfrac{f_k(l)}{l!}  
        \\&= \sum\limits_{l=0}^N \dfrac{f_k(l)}{l!} \sum\limits_{k> N^{\alpha}}\dfrac{1}{k^{1/\alpha}}
        \\ &\le c_2\sum\limits_{l=0}^N \dfrac{1}{(l+1)!} \sum\limits_{k> N^{1/\alpha}}\dfrac{1}{k^{1/\alpha}}
        \\ &\le c_2e\sum\limits_{k> N^{\alpha}}\dfrac{1}{k^{1/\alpha}} 
        = O(N^{-1+\alpha}).
    \end{align*}
    (1) is obtained by substituting $p_k=\dfrac{1}{k^{1/\alpha}\zeta(1/\alpha)}$ and using the inequality $(1 - p_k)^{N-l} \le 1$; (2) follows from the fact that $\zeta(1/\alpha) > 1$ and $k > N^{\alpha}$; (3) uses the bound $\binom{N}{l} \le \dfrac{N^l}{l!}$.

    Therefore, we can see that
    \begin{align*}
   \testdataerror(N)=\mathbb{E}_{k,t_k}
   \left[\sum\limits_{k=1}^{\infty}p_kf_k(t_k) \right]\le O(N^{-1+\alpha}).
    \end{align*}
This completes the proof.   
\end{proof}

\end{document}
